\documentclass[11pt,letterpaper]{article}
\usepackage[margin=1in]{geometry}
\usepackage{graphicx}
\usepackage{booktabs}
\usepackage{longtable}
\usepackage{float}
\usepackage{hyperref}
\usepackage{fancyhdr}
\usepackage{titlesec}

\pagestyle{fancy}
\fancyhead[L]{Gerhard}
\fancyhead[R]{International Analysis}
\fancyfoot[C]{\thepage}

\titleformat{\section}{\Large\bfseries}{\thesection}{1em}{}
\titleformat{\subsection}{\large\bfseries}{\thesubsection}{1em}{}

\begin{document}

\begin{titlepage}
\centering
\vspace*{2cm}
{\Huge\bfseries Gerhard?\par}
\vspace{1cm}
{\LARGE International Tax Revenue Analysis\par}
\vspace{0.5cm}
{\Large Rest of World Comprehensive Report\par}
\vspace{0.5cm}
{\large 1972-2024\par}
\vspace{2cm}
{\large\itshape A Comparative Analysis of Tax Revenue Levels\\
Across 198 Countries\par}
\vfill
{\large Prepared: October 2025\par}
\end{titlepage}

\tableofcontents
\newpage

\section{Executive Summary}

This report analyzes tax revenue levels across 198 countries from 1972-2024, providing a comprehensive view of how different nations finance their governments.

\subsection{Key Findings}

\begin{itemize}
    \item \textbf{Wide Variation}: Tax revenue as a percentage of GDP ranges from near zero to over 44\%, reflecting diverse policy choices and economic structures.
    \item \textbf{Global Median}: The median tax-to-GDP ratio is approximately 15.9\% in recent years.
    \item \textbf{Regional Patterns}: Developed economies typically collect 30-45\% of GDP in taxes, while developing countries often collect 10-20\%.
    \item \textbf{Temporal Trends}: Most developed countries have seen stable or slightly increasing tax levels over the past five decades.
    \item \textbf{Convergence Limited}: Despite globalization, countries maintain distinct tax policy choices with limited convergence.
\end{itemize}

\end{document}
